% ============================================================
% CONTROLE MANUAL DE EXPEDIENTES
% Arquivo: dados/controle-expedientes.tex
%
% Este arquivo NÃO tenta gravar nada automaticamente no disco.
% Ele funciona como um registro simples e editável pelo usuário.
%
% Como usar:
% - Atualize \proximoDiex e \proximoOficio antes de compilar.
% - Em qualquer dados-*.tex de DIEx/Ofício, deixe o número vazio
%   para usar o próximo número informado aqui.
% - Depois de assinar/emitir o expediente, registre-o no histórico
%   e avance manualmente o próximo número.
%
% Exemplo:
%   \diexNumero{}   % usa \proximoDiex
%   \oficioNumero{} % usa \proximoOficio
% ============================================================

\proximoDiex{005}
\proximoOficio{002}

% ------------------------------------------------------------
% Histórico manual de DIEx expedidos
% Formato: \registroDiex{número}{data}{assunto/título}
% ------------------------------------------------------------
\registroDiex{001}{00 de mês de 2026}{notificação prévia}
\registroDiex{002}{00 de mês de 2026}{solicitação de comparecimento de testemunha}
\registroDiex{003}{00 de mês de 2026}{remessa dos autos}
\registroDiex{004}{00 de mês de 2026}{solicitação de prorrogação de prazo}

% ------------------------------------------------------------
% Histórico manual de Ofícios expedidos
% Formato: \registroOficio{número}{data}{assunto/título}
% ------------------------------------------------------------
\registroOficio{001}{00 de mês de 2026}{solicitação de informações}
